% Part: turing-machines
% Chapter: machines-computations
% Section: turing-machines

\documentclass[../../../include/open-logic-section]{subfiles}

\begin{document}

\olfileid{tur}{mac}{tur}
\olsection{ٹیورنگ مشیناں}

\begin{explain}
ٹیورنگ مشین دی
رسمی تعریف ویکھن
نوں تجریدی لگدی اے،
پر اصل وچ سادی اے:
ایہ ٹیورنگ مشین
دا عمل مقرر کرن
لئی درکار ساری
معلومات اک ریاضیاتی
بناوٹ وچ اکٹھیاں
کردی اے۔ ایہناں
وچ (1) مشین دیاں
ممکنہ حالتاں،
(2) پٹی اُتے
آ سکّن والے
نشان، (3) مشین
دی ابتدائی حالت
تے (4) مشین دے
حکماں دا مجموعہ
شامل نیں۔
\end{explain}

\begin{defn}[ٹیورنگ مشین]
اک \emph{ٹیورنگ مشین}
$M$ چار رکناں
دا اک ترتیب وار ٹپل
$\langle Q, \Sigma, q_0,
\delta\rangle$ اے،
جیہدے وچ:
\begin{enumerate}
\item \emph{حالتاں} دا اک
  محدود سیٹ~$Q$،
\item اک محدود
  \emph{حرفاں دا مجموعہ}
  $\Sigma$، جیدے وچ
  $\TMendtape$ تے
  $\TMblank$ شامل نیں،
\item اک \emph{ابتدائی
  حالت}~$q_0 \in Q$،
\item اک محدود
  \emph{حکماں دا مجموعہ}~$\delta\colon Q \times \Sigma
  \pto Q \times \Sigma \times \{\TMleft, \TMright, \TMstay\}$۔
\end{enumerate}
جزوی فنکشن~$\delta$
نوں $M$ دا
\emph{حالتی تبدیلی
دا فنکشن} وی
آکھدے نیں۔
\end{defn}

\begin{explain}
اسی مندے آں کہ
پٹی صرف اک پاسے
ول لامتناہی اے۔
ایس لئی اک خاص
نشان~$\TMendtape$
نوں پٹی دے کھبے
سِرے دی نشانی
بناؤنا فائدہ مند اے۔
ایس نال ٹیورنگ
مشین دے پروگرام
نوں پتا لاؤن وچ
سہولت رہندی اے
کہ اوہ کدوں پٹی
توں باہر نکلن
دے خطرے وچ اے۔
اسی ایہ شرط لا
سکدے ساں کہ ایس
نشان اُتے کدی
نواں نشان نہ لکھیا
جاوے، یعنی
$\delta(q,\TMendtape) = \tuple{q', \TMendtape, x}$
جے $\delta(q,\TMendtape)$
تعریف شدہ ہووے۔
کجھ درسی کتاباں
ایہ شرط رکھدیاں
نیں؛ اسی نہیں
رکھدے۔ ٹیورنگ
مشین بناؤندے
ویلے تسی صرف
ایہ خیال رکھ
سکدے او کہ
اوہ~$\TMendtape$
نوں کدی نہ بدلے۔
ہور ایہ وی
اے کہ کجھ
صورتاں وچ ایس
نشان نوں بدلݨ
دی اجازت دینا
مفید لچک دیندا
اے۔
\end{explain}

\begin{ex}
\emph{جفت مشین:}
جفت مشین رسمی
طور اُتے چار رکناں
دا ترتیب وار ٹپل
$\tuple{Q, \Sigma, q_0, \delta}$
اے، جتھے:
\begin{align*}
Q & = \{ q_0, q_1 \} \\
\Sigma & = \{ \TMendtape, \TMblank, \TMstroke \}, \\
\delta(q_0, \TMstroke) & = \tuple{q_1, \TMstroke, \TMright},\\
\delta(q_1, \TMstroke) & = \tuple{q_0, \TMstroke, \TMright},\\
\delta(q_1, \TMblank)  & = \tuple{q_1, \TMblank, \TMright}.
\end{align*}
\end{ex}

\end{document}
